% SPDX-FileCopyrightText: 2026 Will Cook
% SPDX-License-Identifier: CC-BY-4.0
%
% Expository paper for Erdős Problem #257: the theorem of Appendix A of the
% short paper, explained with examples in Sections 1 to 8 and proved in full in
% Appendix A. It goes with the films that explain the same argument.
\documentclass[10pt]{article}
\usepackage{paper-house-style}
\usepackage{tikz}
\usetikzlibrary{decorations.pathreplacing,calc}
\usepackage{xurl}
\usepackage[colorlinks=true,linkcolor=linkinternal,urlcolor=linkexternal,citecolor=linkinternal]{hyperref}
\hypersetup{pdftitle={Sums of 1/(2ⁿ−1) that are never fractions},pdfauthor={Will Cook},pdflang={en-GB},bookmarksnumbered=true}
\runninghead{Sums of $1/(2^n-1)$ that are never fractions}
\title{Sums of $1/(2^n-1)$ that are never fractions}
\subtitle{An irrationality theorem stated by Erd\H{o}s, explained with examples\\ and proved in full (Erd\H{o}s Problem \#257)}
\author{Will Cook}
\date{11 October 2026}
\emergencystretch=1.5em
\renewcommand{\bottomfraction}{0.4}

\theoremstyle{plain}
\newtheorem{theorem}{Theorem}
\newtheorem{lemma}[theorem]{Lemma}
\newtheorem{corollary}[theorem]{Corollary}
\theoremstyle{definition}
\newtheorem{example}[theorem]{Example}
\newtheorem{problem}{Problem}
\newtheorem{exercise}{Exercise}
\newtheorem{readercheck}{Question}
\theoremstyle{remark}
\newtheorem*{remark}{Remark}

\newcommand{\lft}{\ell}
\newcommand{\rmod}{\mskip5mu\mathrm{mod}\mskip5mu}

\begin{document}
\maketitle
\paperprovenance

\begin{abstract}
Take the powers of $2$, subtract $1$ from each, and put each result under a
$1$: this gives $1,\tfrac13,\tfrac17,\tfrac1{15},\tfrac1{31},\dots$. Erd\H{o}s
asked whether adding up an infinite subset of this list, that is, infinitely
many of these fractions with each used at most once, can ever give a
fraction. Each fraction comes from an exponent $n$: the exponent $3$ gives
$\tfrac1{2^3-1}=\tfrac17$. We prove that the chosen fractions add up to an
irrational number whenever $1/n$, added up over the chosen exponents $n$, is
finite. For example, for the exponents $2,4,8,16,\dots$ the sum of $1/n$ is $\tfrac12+\tfrac14+\tfrac18+\dots$, which adds up to $1$
(Section~1 shows why), so the chosen fractions
$\tfrac13+\tfrac1{15}+\tfrac1{255}+\tfrac1{65535}+\dots$ add up to an
irrational number. Sections~1 to~8 explain the proof with examples, and
Appendix~A gives every step in full. A computer has checked a proof of this
theorem in the Lean proof assistant; the version written here is checked by
hand. For an arbitrary infinite choice the question is still open.
\end{abstract}

\section{What Erd\H{o}s asked}\label{sec:intro}

Start with the powers of $2$, subtract $1$ from each, and put each result
under a $1$:
\begin{center}
\begin{tabular}{l|ccccccc}
\papertablehead
exponent $n$ & $1$ & $2$ & $3$ & $4$ & $5$ & $6$ & $\dots$\\
$2^n$ & $2$ & $4$ & $8$ & $16$ & $32$ & $64$ & $\dots$\\
$2^n-1$ & $1$ & $3$ & $7$ & $15$ & $31$ & $63$ & $\dots$\\
fraction $\dfrac1{2^n-1}$ & $1$ & $\tfrac13$ & $\tfrac17$ & $\tfrac1{15}$ & $\tfrac1{31}$ & $\tfrac1{63}$ & $\dots$
\end{tabular}
\end{center}
The top row is the \emph{exponent} $n$, and the exponent $n$ gives the
fraction $1/(2^n-1)$: the exponent $3$ gives $\tfrac17$, because
$2^3-1=8-1=7$. The exponents are $1,2,3,\dots$; the exponent $0$ is left out,
since $2^0-1=0$.

Now choose an infinite subset of the list: infinitely many of its fractions,
each used at most once. Choosing the subset is the same as choosing its
exponents. Add up the fractions in it, for example:
\begin{center}
\begin{tabular}{ll}
every exponent $1,2,3,4,\dots$: & $1+\tfrac13+\tfrac17+\tfrac1{15}+\dots$\\[3pt]
the even exponents $2,4,6,8,\dots$: & $\tfrac13+\tfrac1{15}+\tfrac1{63}+\tfrac1{255}+\dots$\\[3pt]
the doubling exponents $2,4,8,16,\dots$: & $\tfrac13+\tfrac1{15}+\tfrac1{255}+\tfrac1{65535}+\dots$
\end{tabular}
\end{center}
Erd\H{o}s asked whether any choice like this can add up to a fraction:
exactly $\tfrac25$, say, or $\tfrac{17}{41}$, or any other. A number that
cannot be written as a fraction is called \emph{irrational}. Its decimal never
ends and never settles into a repeating block, like $\sqrt2=1.41421356\dots$
(Exercise~\ref{ex:5} proves that $\sqrt2$ is irrational), while a fraction's decimal either ends, like $\tfrac34=0.75$, or repeats for
ever, like $\tfrac17=0.142857\,142857\dots$.

We can never finish adding infinitely many numbers, so we add them one at a
time and watch the \emph{running total}. Take
\[
 \frac12+\frac14+\frac18+\frac1{16}+\dots .
\]
The running totals are $\tfrac12,\tfrac34,\tfrac78,\tfrac{15}{16},\dots$, and
the gaps to $1$ are $\tfrac12,\tfrac14,\tfrac18,\tfrac1{16},\dots$. Each new
number is exactly half of the gap left before it: $\tfrac14$ is half of the
gap $\tfrac12$, and $\tfrac18$ is half of the gap $\tfrac14$. So the gap
halves every time (Figure~\ref{fig:halving}). The running totals therefore
never reach $1$, and they get as close to $1$ as we like: after $10$ numbers
the gap is $\tfrac1{1024}$, and after $20$ it is $\tfrac1{1048576}$, less than
a millionth. That is what
\[
 \frac12+\frac14+\frac18+\frac1{16}+\dots=1
\]
means. So infinitely many fractions \emph{can} add up to a fraction, and a
proof that our sums are irrational has to use something special about the
bottoms $1,3,7,15,\dots$.

\begin{figure}[htb]
\centering
\begin{tikzpicture}[x=7.2cm,y=0.80cm]
 \node[anchor=east,font=\small] at (-0.06,0.65) {numbers added};
 \node[anchor=west,font=\small] at (1.08,0.65) {gap to $1$};
 \foreach \v/\yy/\termsadded/\total/\gap in
   {0.5/0/1/\tfrac12/\tfrac12,
    0.75/-1/2/\tfrac34/\tfrac14,
    0.875/-2/3/\tfrac78/\tfrac18,
    0.9375/-3/4/\tfrac{15}{16}/\tfrac1{16}}
 {
  \fill[gray!22] (0,\yy-0.10) rectangle (\v,\yy+0.10);
  \draw (0,\yy) -- (\v,\yy);
  \draw[densely dashed,thick] (\v,\yy) -- (1,\yy);
  \draw (0,\yy-0.14) -- (0,\yy+0.14);
  \draw (1,\yy-0.14) -- (1,\yy+0.14);
  \fill (\v,\yy) circle (1.6pt);
  \node[above=2pt,font=\small] at (\v,\yy) {$\total$};
  \node[anchor=east,font=\small] at (-0.06,\yy) {$\termsadded$};
  \node[anchor=west,font=\small] at (1.08,\yy) {$\gap$};
 }
 \node[font=\small] at (0,-3.50) {$0$};
 \node[font=\small] at (1,-3.50) {$1$};
 \node[font=\small] at (-0.06,-3.95) {$\vdots$};
 \node[font=\small] at (1.08,-3.95) {$\vdots$};
 \node[font=\small] at (0.50,-4.10) {value of the running total};
\end{tikzpicture}
\caption{The running totals of $\tfrac12+\tfrac14+\tfrac18+\dots$, each on a
line from $0$ to $1$; the dashed part is the gap to $1$, and each new number
fills half of it. The rows go on for ever.}
\label{fig:halving}
\end{figure}

Erd\H{o}s proved in 1948 that choosing every exponent gives an irrational sum
\cite{erdos1948}. Whether every infinite choice does is Erd\H{o}s Problem
\#257 \cite{erdosproblems257}, and it is still open.

\section{What this paper proves}\label{sec:problem}

We write $A$ for the set of chosen exponents, which is an infinite subset of
$1,2,3,\dots$, and $\sum_{n\in A}$ for ``add up
over the exponents $n$ in $A$''.

\begin{theorem}\label{thm:main}
Let $A$ be an infinite set of positive integers such that the sum
$\sum_{n\in A}1/n$ is finite. Then
\[
 \sum_{n\in A}\frac1{2^n-1}
\]
is irrational.
\end{theorem}

For $A=\{2,4,8,16,\dots\}$ the theorem reads: $A$ is infinite;
$\sum_{n\in A}1/n=\tfrac12+\tfrac14+\tfrac18+\dots=1$ is finite; so
$\sum_{n\in A}1/(2^n-1)=\tfrac13+\tfrac1{15}+\tfrac1{255}+\dots$ is
irrational. A sum of infinitely many positive numbers is \emph{finite} when
its running totals get as close as we like to one number without passing it,
as the running totals of $\tfrac12+\tfrac14+\dots$ do with $1$;
Appendix~\ref{app:sums} says exactly what that means and shows that every
choice of exponents gives a definite sum. Two more examples:

\begin{itemize}
\item \emph{The squares} $1,4,9,16,\dots$. Each $\tfrac1{m^2}$ after the first
 is less than $\tfrac1{(m-1)m}$, which equals $\tfrac1{m-1}-\tfrac1m$: for
 instance $\tfrac14<\tfrac12=1-\tfrac12$ and $\tfrac19<\tfrac16=\tfrac12-\tfrac13$.
 These differences cancel in a chain,
 \[
  \Bigl(1-\frac12\Bigr)+\Bigl(\frac12-\frac13\Bigr)+\dots+\Bigl(\frac1{m-1}-\frac1m\Bigr)=1-\frac1m ,
 \]
 so $\tfrac14+\tfrac19+\dots+\tfrac1{m^2}$ stays below $1$ however far we
 go, and $1+\tfrac14+\tfrac19+\dots$ is finite: its running totals only go up
 and stay below $2$ (Appendix~\ref{app:sums}). The theorem says that
 $1+\tfrac1{15}+\tfrac1{511}+\tfrac1{65535}+\dots$ is irrational.
\item \emph{Every exponent} $1,2,3,4,\dots$. Now $1+\tfrac12+\tfrac13+\tfrac14+\dots$
 is infinite: $\tfrac13+\tfrac14$ is at least $\tfrac14+\tfrac14=\tfrac12$;
 the next four numbers, $\tfrac15$ to $\tfrac18$, are each at least $\tfrac18$,
 so they add up to at least $\tfrac12$; the next eight add up to at least
 $\tfrac12$ again; and so on, so the running totals pass every number. The
 theorem says nothing about this choice. (Erd\H{o}s's 1948 proof covers it by
 another method.)
\end{itemize}
The theorem also needs $A$ to be infinite: for $A=\{2,4\}$,
$\tfrac12+\tfrac14=\tfrac34$ is finite, but $\tfrac13+\tfrac1{15}=\tfrac25$ is
a fraction.

In 1968 Erd\H{o}s proved Theorem~\ref{thm:main} when no two chosen exponents
have a common factor bigger than $1$, and wrote that this condition could be
removed, without giving the details \cite[p.~222]{erdos1968}. Tao and
Ter\"av\"ainen settled the prime exponents $2,3,5,7,11,\dots$
\cite[Theorem~1.3]{taoteravainen2025}; for the primes $\sum_{n\in A}1/n$ is
infinite (a theorem of Euler from 1737), so Theorem~\ref{thm:main} does not
cover them. Theorem~\ref{thm:main}
with any whole number $b\ge2$ in place of $2$ is Theorem~A.1 of the short
paper \cite{shortpaper}, and a computer has checked a proof of it in the Lean
proof assistant (Section~\ref{sec:limits}). This paper explains the case
$b=2$ for a reader new to the problem. It follows
\cite[Appendix~A]{shortpaper} and \cite[Section~1.1]{longrecord} with one
change in the last step, explained in Appendix~\ref{app:lean}. That version is
checked by hand, not by Lean.

\section{The test: a fraction cannot leave a tiny positive remainder}\label{sec:test}

Suppose the sum $S$ of the chosen fractions were a fraction $p/q$, with $p$
and $q$ whole numbers.
Multiply it by a whole number $m$ and take away a whole number $j$:
\[
 mS-j=\frac{mp}q-\frac{jq}q=\frac{mp-jq}q .
\]
The top, $mp-jq$, is one whole number minus another. So if $mS-j$ is
positive, its top is a positive whole number, which is at least $1$, and
$mS-j$ is at least $\tfrac1q$.

For example, if $S$ were $\tfrac25$:
\[
 3\times\frac25-1=\frac15,\qquad 4\times\frac25-1=\frac35,\qquad
 5\times\frac25-2=0,\qquad 1\times\frac25-1=-\frac35 .
\]
Every answer lands on one of the marks
$\dots,-\tfrac25,-\tfrac15,0,\tfrac15,\tfrac25,\dots$ (Figure~\ref{fig:fifths}),
and there is no mark between $0$ and $\tfrac15$. So a positive answer is at
least $\tfrac15$, and no choice of $m$ and $j$ gives $0.00001$.
The word \emph{positive} matters, since an answer can also be $0$ or negative.

\begin{figure}[htb]
\centering
\begin{tikzpicture}[x=6.5cm,y=1cm]
 \fill[gray!20] (0,-0.10) rectangle (0.2,0.10);
 \draw[thick,<->] (-0.53,0) -- (1.33,0);
 \foreach \v/\lab in {-0.4/-\tfrac25,-0.2/-\tfrac15,0/0,
    0.2/\tfrac15,0.4/\tfrac25,0.6/\tfrac35,0.8/\tfrac45,1/1,1.2/\tfrac65}
 {
  \draw (\v,0.12) -- (\v,-0.12);
  \fill (\v,0) circle (1.5pt);
  \node[below=4pt,font=\small] at (\v,0) {$\lab$};
 }
 \draw[->] (0.1,0.55) -- (0.1,0.15);
 \node[above,font=\small] at (0.1,0.55) {no mark strictly inside this gap};
 \node[font=\small] at (0.4,-0.95) {values of $m\times\tfrac25-j$ for whole numbers $m$ and $j$};
\end{tikzpicture}
\caption{Part of the line of fifths, which continues in both directions. Every value of $m\times\tfrac25-j$ lies on a mark, and the shaded
gap between $0$ and $\tfrac15$ contains no mark.}
\label{fig:fifths}
\end{figure}

This gives the strategy of the whole proof. Whatever bottom $q$ someone
claims for $S$, find whole numbers $m$ and $j$ with
\[
 0<mS-j<\frac1q .
\]
Then $S$ is not $p/q$, and doing this for every $q$ shows that $S$ is not a
fraction at all. The numbers $m$ and $j$ are allowed to depend on $q$.

\section{The doubling exponents: multiply by \texorpdfstring{$2^k-1$}{2\^{}k-1}}\label{sec:doubling}

Take $S=\tfrac13+\tfrac1{15}+\tfrac1{255}+\tfrac1{65535}+\dots$ and multiply
it by $255=2^8-1$. The first three fractions give whole numbers, and the rest
give very small ones:
\[
 \frac{255}3=85,\quad \frac{255}{15}=17,\quad \frac{255}{255}=1,\quad
 \frac{255}{65535}=\frac1{257},\quad \frac{255}{2^{32}-1}=\frac1{16843009},\quad\dots .
\]
So, with $103=85+17+1$,
\[
 255S=103+\Bigl(\frac1{257}+\frac1{16843009}+\dots\Bigr).
\]
Take away $103$. What is left,
\[
 R(8)=255S-103=\frac1{257}+\frac1{16843009}+\dots\approx0.0039 ,
\]
is positive and small. This is the test with $m=255$ and $j=103$. If $S$
were $p/q$, then $R(8)$ would be at least $\tfrac1q$; since $R(8)$ is less
than $\tfrac1{64}$ (shown below), no bottom $q$ from $1$ to $64$ is possible.

The multiplier $255$ works because $3$, $15$ and $255$ all divide it.
Section~\ref{sec:split} shows that this happens because the exponents $2$,
$4$ and $8$ all divide $8$. Multiplying by $2^{16}-1=65535$
instead makes the first four fractions whole numbers; the exponent $32$ leaves
$\tfrac{65535}{2^{32}-1}=\tfrac1{65537}$, and what is left in total is about
$0.000015$.

\paragraph{Why this works for every bottom.} Let $k$ be any of the doubling
exponents, multiply $S$ by $2^k-1$, and look at each chosen fraction.
\begin{itemize}
\item Every chosen exponent up to $k$ divides $k$, so those fractions give
 whole numbers.
\item Every later chosen exponent $n$ is at least $2k$, and its fraction gives
 $\frac{2^k-1}{2^n-1}$. This is less than $\frac{2^k}{2^{n-1}}=\frac2{2^{n-k}}$,
 because $2^n-1$ is at least half of $2^n$.
\end{itemize}
Take away the whole numbers, and call what is left $R(k)$; it is positive,
because there is always a chosen exponent beyond $k$. The gaps $n-k$ between
$k$ and the later exponents are different whole numbers, all at least $k$
(for $k=8$ they are $8,24,56,\dots$). Allowing every gap from $k$ upwards only
adds more positive numbers, so
\[
 0<R(k)<\frac2{2^k}+\frac2{2^{k+1}}+\frac2{2^{k+2}}+\dots=\frac4{2^k}.
\]
The last step uses the same pattern as $\tfrac12+\tfrac14+\dots=1$: when each
number is half the one before, the sum is twice the first number. For $k=8$
the bound is $\tfrac1{64}$, and for $k=16$ it is $\tfrac1{16384}$.

Now take any bottom $q$, and choose a doubling exponent $k$ with $2^k>4q$.
Then $R(k)<\tfrac4{2^k}<\tfrac1q$. But $R(k)$ is $S$ times the whole number
$2^k-1$, minus a whole number, and it is positive, so if $S$ were $p/q$ it
would be at least $\tfrac1q$. It cannot be both. So
$\tfrac13+\tfrac1{15}+\tfrac1{255}+\dots$ is irrational. That is the whole
proof for the doubling exponents (Appendix~\ref{app:doubling} writes it out as
Theorem~\ref{thm:doubling}).

\section{Why \texorpdfstring{$2^n-1$}{2\^{}n-1} divides \texorpdfstring{$2^k-1$}{2\^{}k-1}, and what is left when it does not}\label{sec:split}

A whole number divided by another splits into a \emph{whole part} and a
\emph{fractional part} that is at least $0$ and less than $1$: for example
$\tfrac{255}7=36+\tfrac37$, with whole part $36$ and fractional part $\tfrac37$. For
our numbers the fractional part follows one rule.

Going from $k$ to $k+n$ multiplies $2^k$ by $2^n$, and $2^n$ is one more than
$2^n-1$. So $2^{k+n}-1=2^k(2^n-1)+(2^k-1)$, and dividing by $2^n-1$,
\[
 \frac{2^{k+n}-1}{2^n-1}=2^k+\frac{2^k-1}{2^n-1}.
\]
With $n=3$, for example, $\tfrac{63}7=8+\tfrac77$ (Figure~\ref{fig:dots}) and
$\tfrac{255}7=32+\tfrac{31}7$. So adding $n$ to $k$ adds the whole number
$2^k$, and adding a whole number does not change the fractional part. For the exponent $n$ the fractional
part therefore repeats every $n$ steps of $k$, and depends only on the
remainder $r$ when $k$ is divided by $n$. At $k=r$ the fraction
$\tfrac{2^r-1}{2^n-1}$ is less than $1$, because $r$ is less than $n$, so it
is its own fractional part. Hence
\[
 \text{the fractional part of }\frac{2^k-1}{2^n-1}\text{ is }\frac{2^r-1}{2^n-1},
 \qquad\text{where $r$ is the remainder when $k$ is divided by $n$.}
\]
For $k=8$ and $n=3$: $8=2\times3+2$, so $r=2$, and the fractional part is
$\tfrac{2^2-1}7=\tfrac37$, as above.

The fractional part is $0$ exactly when $r=0$, that is, when $n$ divides $k$; we
then say that $k$ \emph{cancels} the exponent $n$. So $2^n-1$ divides $2^k-1$
exactly when $n$ divides $k$: $15=2^4-1$ divides
$255=2^8-1$ because $4$ divides $8$, while $7=2^3-1$ does not, because $3$
does not divide $8$. Here is the exponent $3$, with the fraction $\tfrac17$, as $k$
goes up:
\begin{center}
\begin{tabular}{l|cccccccccc}
\papertablehead
$k$ & $1$ & $2$ & $3$ & $4$ & $5$ & $6$ & $7$ & $8$ & $9$ & $\dots$\\
remainder of $k$ divided by $3$ & $1$ & $2$ & $0$ & $1$ & $2$ & $0$ & $1$ & $2$ & $0$ & $\dots$\\
fractional part of $\tfrac{2^k-1}7$ & $\tfrac17$ & $\tfrac37$ & $0$ & $\tfrac17$ & $\tfrac37$ & $0$ & $\tfrac17$ & $\tfrac37$ & $0$ & $\dots$
\end{tabular}
\end{center}

\begin{figure}[htb]
\centering
\begin{tikzpicture}[scale=0.38]
 \foreach \r in {0,...,6}{\foreach \c in {0,...,7}{\fill (\c,\r) circle (0.16);}}
 \foreach \c in {0,...,6}{\fill (\c,7) circle (0.16);}
 \draw (7,7) circle (0.16);
 \foreach \r in {0,...,7}{\draw[rounded corners=2pt,gray] (-0.4,\r-0.4) rectangle (6.4,\r+0.4);}
 \draw[rounded corners=2pt,gray,thick] (6.6,-0.4) rectangle (7.4,6.4);
 \node[anchor=west,font=\small,align=left] at (8.2,5.2) {$2^6=8\times8=64$ dots;\\ remove one: $63$};
 \node[anchor=west,font=\small,align=left] at (8.2,2.2) {$8$ rows of $7$, and\\ one column of $7$:\\ $63=8\times7+7=9\times7$};
\end{tikzpicture}
\caption{The step from $k=3$ to $k=6$ for the exponent $3$: $2^6-1=2^3\times7+7$,
so $\tfrac{63}7=8+\tfrac77=9$. The open circle is the dot removed.}
\label{fig:dots}
\end{figure}

\paragraph{The sum of the fractional parts.} Now take any infinite set of exponents $A$,
with sum $S$, and any whole number $k\ge1$. Multiply every chosen fraction by
$2^k-1$, take away the whole part of each, and add up the fractional parts. Call the
result $R(k)$. Two facts make it useful.
\begin{itemize}
\item Only exponents up to $k$ can have a whole part other than $0$, because
 a later fraction times $2^k-1$ is less than $1$. So the whole parts add up to a whole
 number $J(k)$, and $R(k)=(2^k-1)S-J(k)$: this is the test of
 Section~\ref{sec:test} with $m=2^k-1$ and $j=J(k)$.
\item $R(k)$ is positive, because $A$ is infinite, so it has a chosen exponent
 beyond $k$, and that fraction times $2^k-1$ is all fractional part.
\end{itemize}
So if $S$ were $p/q$, then \emph{every} $R(k)$ would be at least $\tfrac1q$
(Lemma~\ref{lem:split} and Corollary~\ref{cor:atleast}). The rest of the proof
finds one $k$ with $R(k)<\tfrac1q$. ($R(k)$ adds the separate fractional
parts, so it can be bigger than $1$; Appendix~\ref{app:split} has an
example.)

\section{Other choices: averaging}\label{sec:general}

\paragraph{Why cancelling the early exponents is not enough.} For the squares
$1,4,9,16,\dots$, both $k=4$ and $k=8$ cancel exponents $1$ and $4$, which
divide both. Every other square is at least $9$, bigger than $8$, so each of
those fractions is all fractional part, and $R(k)$ is the multiplier
times the same sum $\tfrac1{511}+\tfrac1{65535}+\tfrac1{2^{25}-1}+\dots$ in
both cases. The multipliers are $15$ and $255=17\times15$, so
\[
 R(8)=17\times R(4)\approx0.503,\qquad R(4)\approx0.0296 .
\]
The bigger $k$ kept the same exponents cancelled and made everything else
seventeen times bigger; the exponent $9$ alone leaves
$\tfrac{255}{511}\approx0.499$. Cancelling longer batches does not fix this
by itself: Exercise~\ref{ex:6} gives a choice of exponents and, for every
$j$, a $k$ that cancels the first $j$ of them while the very next chosen
exponent, $k+1$, leaves at least $\tfrac37$.

So we use many values of $k$ at once. Fix a whole number $Q$, the
\emph{step}, that every exponent in an early batch of chosen exponents
divides, and take
\[
 k=Q,\ 2Q,\ 3Q,\ \dots,\ MQ .
\]
Every one of these cancels the batch. If the values
$R(Q),R(2Q),\dots,R(MQ)$ have an average below $\tfrac1q$, then at least one
of them is below $\tfrac1q$, because if every one were at least $\tfrac1q$, so
would be their average. (For example, $3$, $5$ and $10$ have average $6$, and
$3$ is below $6$.)

\paragraph{One exponent over a full cycle.} The exponent $3$ leaves $\tfrac17$,
$\tfrac37$, $0$ and then repeats. Over one full cycle these average
\[
 \frac13\Bigl(\frac17+\frac37+0\Bigr)=\frac4{21}\approx0.19,
\]
less than $\tfrac13$. Here is why. Add $\tfrac17$ to each fractional part to get
$\tfrac27,\tfrac47,\tfrac17$. These add up to exactly $1$, since $2+4+1=7$, so
they average exactly $\tfrac13$, and the fractional parts themselves average
$\tfrac17$ less (Figure~\ref{fig:bars}). One of them, $\tfrac37$, is bigger than
$\tfrac13$; it is the average that is below $\tfrac13$. The same works for every
exponent $n$, because $1+2+4+\dots+2^{n-1}=2^n-1$: over a full cycle the
fractional parts for the exponent $n$ average $\tfrac1n-\tfrac1{2^n-1}$, less
than $\tfrac1n$.

\begin{figure}[htb]
\centering
\begin{tikzpicture}[x=0.78cm,y=0.50cm]
 \draw[->] (0,0) -- (7.6,0) node[right,font=\small]{$k$};
 \draw[->] (0,0) -- (0,5.2);
 \foreach \y/\l in {0/0,1/\tfrac17,2/\tfrac27,3/\tfrac37,4/\tfrac47}
   {\draw (-0.1,\y) -- (0.1,\y) node[left=3pt,font=\scriptsize]{$\l$};}
 \foreach \k/\h in {1/1,2/3,3/0,4/1,5/3,6/0}
   {\fill[gray!55] (\k-0.3,0) rectangle (\k+0.3,\h);
    \draw[gray!70!black] (\k-0.3,\h) rectangle (\k+0.3,\h+1);
    \fill[gray!12] (\k-0.3+0.02,\h+0.04) rectangle (\k+0.3-0.02,\h+1-0.04);
    \node[below,font=\scriptsize] at (\k,0) {\k};}
 \draw[densely dashed,thick] (0,1.3333) -- (6.6,1.3333);
 \node[right,font=\scriptsize] at (6.6,1.3333) {average of the fractional parts, $\tfrac4{21}$};
 \draw[thick] (0,2.3333) -- (6.6,2.3333);
 \node[right,font=\scriptsize] at (6.6,2.3333) {average with $\tfrac17$ added, $\tfrac13$};
 \node[font=\scriptsize] at (7.2,-0.55) {$\dots$};
 \node[font=\small,anchor=west] at (0,-1.9) {fractional part of $\tfrac{2^k-1}7$ for each $k$ (dark), with $\tfrac17$ added (light)};
\end{tikzpicture}
\caption{The exponent $3$ leaves $\tfrac17,\tfrac37,0$ and repeats (dark bars).
With $\tfrac17$ added (light caps) they are $\tfrac27,\tfrac47,\tfrac17$, which
average exactly $\tfrac13$ (solid line), so the fractional parts average
$\tfrac13-\tfrac17=\tfrac4{21}$ (dashed line). Only the averages are compared:
the bar $\tfrac37$ is above $\tfrac13$.}
\label{fig:bars}
\end{figure}

\paragraph{Steps of $Q$.} With $k$ going up in steps of $Q$, an exponent can
skip some remainders. Take the exponent $6$ with $Q=4$: the values $k=4,8,12$
leave remainders $4,2,0$ when divided by $6$, then repeat, so the fractional parts
are $\tfrac{15}{63},\tfrac3{63},0$, averaging $\tfrac2{21}$, which is less than
$\tfrac16$. The skipped remainders do not spoil the bound, for a reason:
the remainders visited are exactly $0,2,4$, the multiples of $2$, the highest
common factor of $6$ and $4$, each once per cycle. With $\tfrac1{63}$ added,
the fractional parts are $\tfrac{16}{63},\tfrac4{63},\tfrac1{63}$, which add up to
$\tfrac{21}{63}=\tfrac13$; so they average $\tfrac19$, and the fractional parts
average $\tfrac19-\tfrac1{63}=\tfrac2{21}$. In general, over a full cycle of $k=Q,2Q,3Q,\dots$,
exponent $n$ averages less than $\tfrac1n$ (Lemma~\ref{lem:cycle}).

\paragraph{Unfinished cycles.} Over part of a cycle the average can be
bigger than $\tfrac1n$. The exponent $9$ with $Q=4$ leaves $\tfrac{15}{511},\tfrac{255}{511},\tfrac7{511},\tfrac{127}{511}$
at the first four values of $k$, which average $\tfrac{101}{511}$, more than
$\tfrac19$; the full cycle of nine values averages $\tfrac19-\tfrac1{511}$,
below $\tfrac19$. With infinitely many exponents, a fixed number $M$ of values
of $k$ always leaves some exponents unfinished: a chosen exponent $n$ bigger
than $MQ$ never even reaches remainder $0$. For those exponents we use a
weaker bound that holds however few values of $k$ we take:
\[
 \text{the average fractional part for the exponent $n$ over }k=Q,2Q,\dots,MQ
 \text{ is at most }\frac{Q+2}n .
\]
The idea: a fractional part for the exponent $n$ is large only when $k$ is just below a
multiple of $n$. If $s$ is the distance from $k$ up to the next multiple of
$n$, the fractional part is less than $\tfrac1{2^s}+\tfrac1{2^{s+n}}+\tfrac1{2^{s+2n}}+\dots$.
As $k$ runs through $Q,2Q,\dots,MQ$, each distance $s$ comes up at most
$\tfrac{MQ+s}n$ times, and adding up gives at most $\tfrac{MQ+2}n$ in total
(Lemma~\ref{lem:crude}). The bound grows with $Q$ but not with $M$, so taking
more values of $k$ never spoils it.

\section{The proof of Theorem~\ref{thm:main}}\label{sec:proof}

Suppose the sum $S$ were $p/q$. We find values of $k$ whose values $R(k)$
average less than $\tfrac1q$. Picture a table with a row for each chosen
exponent and a column for each $k=Q,2Q,\dots,MQ$; each entry is the
fractional part for that exponent at that $k$, and each column adds up to $R(k)$. We call
this table the \emph{grid}, and $R(k)$ its bottom row. The
average of the column totals equals the sum of the row averages, since both
add up the same entries (Appendix~\ref{app:proof} checks this for infinitely
many rows). Split the rows at two cutoffs $t\le T$, chosen below, into three
groups:
\begin{center}
\small
\begin{tabular}{>{\raggedright\arraybackslash}p{2.2cm}|>{\raggedright\arraybackslash}p{5.2cm}|>{\raggedright\arraybackslash}p{4.4cm}}
\papertablehead
exponents & what we do & contribution to the average\\
early, $n\le t$ & choose $Q$ so that every one of them divides $Q$ & $0$\\
middle, $t<n\le T$ & choose $M$ so that every one of them finishes whole cycles & less than $\sum_{n\in A,\,n>t}\tfrac1n$\\
far, $n>T$ & use the bound $\tfrac{Q+2}n$ & at most $(Q+2)\sum_{n\in A,\,n>T}\tfrac1n$
\end{tabular}
\end{center}
The four numbers are chosen in this order, each with a job.
\begin{enumerate}
\item[$t$:] Since $\sum_{n\in A}1/n$ is finite, the part of it beyond a
 cutoff gets as small as we like as the cutoff grows. Choose $t$ so that the
 part beyond $t$ is less than $\tfrac1{2q}$.
\item[$Q$:] Multiply together the chosen exponents up to $t$ (or take $Q=1$ if
 there are none). Every one of them divides $Q$, so they leave $0$ at every
 $k=Q,2Q,\dots$.
\item[$T$:] Now that $Q$ is fixed, choose $T$, at least $t$, so that $Q+2$
 times the part beyond $T$ is less than $\tfrac1{2q}$. This is why $T$ comes
 after $Q$: the far bound contains $Q+2$.
\item[$M$:] There are only finitely many middle exponents. Take $M$ to be the
 product of their cycle lengths (or $1$ if there are none), so that every
 middle row finishes whole cycles.
\end{enumerate}
Then the early rows contribute $0$, the middle rows less than $\tfrac1{2q}$,
and the far rows less than $\tfrac1{2q}$. So
\[
 \frac{R(Q)+R(2Q)+\dots+R(MQ)}M<\frac1{2q}+\frac1{2q}=\frac1q ,
\]
and at least one of these values $R(k)$ is less than $\tfrac1q$. But every
$R(k)$ is positive, so if $S$ were $p/q$, every one would be at least
$\tfrac1q$ (Section~\ref{sec:split}). Both cannot hold. The same argument
works whatever bottom $q$ we start from, so $S$ is not a fraction. That proves
Theorem~\ref{thm:main}.

Each assumption did one job. $A$ being infinite made every $R(k)$
positive. Supposing $S=p/q$ made every $R(k)$ at least $\tfrac1q$. And
$\sum_{n\in A}1/n$ being finite made the parts beyond $t$ and beyond $T$
small. Appendix~\ref{app:numbers} runs the four choices with actual numbers:
for the doubling exponents and $q=5$ they are $t=16$, $Q=1024$, $T=16384$ and
$M=1024$.

\section{What the proof shows, and what it does not}\label{sec:limits}

\paragraph{Every exponent.} If every exponent is chosen, then for any $k$ the
exponents beyond $k$ are all fractional part, and
\[
 R(k)>\frac{2^k-1}{2^{k+1}}+\frac{2^k-1}{2^{k+2}}+\dots=\frac{2^k-1}{2^k}\ge\frac12 ,
\]
since replacing each bottom $2^n-1$ by the bigger $2^n$ makes each fraction
smaller, and the halving list on the right adds up to twice its first number.
So $R(k)$ never gets below $\tfrac12$ (for instance
$R(7)\approx2.05$), and this proof cannot reach this choice. Its sum is still
irrational, by Erd\H{o}s's 1948 argument \cite{erdos1948}. A proof for every
infinite set of exponents will need more than small values of $R(k)$.

\paragraph{Other bases.} The same proof works with any whole number $b\ge2$ in
place of $2$, giving the sums of $1/(b^n-1)$, for example
$\tfrac18+\tfrac1{80}+\tfrac1{6560}+\dots$ for $b=3$ and exponents
$2,4,8,\dots$ (Appendix~\ref{app:bases}).

\paragraph{What is open.} Theorem~\ref{thm:main} answers Erd\H{o}s's question
for every choice with $\sum_{n\in A}1/n$ finite. The short paper
\cite{shortpaper} proves irrationality under weaker conditions, which include
some choices where $\sum_{n\in A}1/n$ is infinite, and the long record
\cite{longrecord} keeps every route tried on the full problem, including those
that failed and why. For an arbitrary infinite choice the question is open.

\paragraph{The computer-checked proof.} Lean is a programming language in which
a computer checks every step of a proof. The theorem for every base $b\ge2$ is
the Lean theorem \texttt{irrational\_erdosSupportSeries\_of\_summable\_reciprocal}
(\href{https://github.com/wcook04/plectis-erdos/blob/436f55ebdafa67e4af0fff79f621c13f2ded12bf/lean/Erdos249257/AllBaseReciprocalSupportIrrationality.lean#L395}{source at a fixed commit}).
Its proof takes about two thousand lines in its own file and the file it
builds on, and reading it takes fluency in Lean. A separate tool, Comparator,
checked that the proof proves exactly a separately written copy of the
statement, using only Lean's three standard axioms (replay run 35935225572).
Neither check says that the Lean statement means the sentence of
Theorem~\ref{thm:main}; Appendix~\ref{app:lean} reads it line by line so the
reader can compare. The argument of Sections~\ref{sec:test}
to~\ref{sec:proof} is this paper's rewriting, checked by hand.

\appendix

\section{The proof in full}\label{app:full}

This appendix proves every step of Sections~\ref{sec:test}
to~\ref{sec:proof}, the facts about infinite sums used in
Sections~\ref{sec:intro} and~\ref{sec:problem}, and the other bases of
Section~\ref{sec:limits}. Whole numbers are $0,1,2,3,\dots$, and positive integers
are $1,2,3,\dots$.

\subsection{Adding infinitely many numbers}\label{app:sums}

\paragraph{Adding up to a number.} Infinitely many positive numbers \emph{add
up to} a number $L$ when their running totals never pass $L$ and get as close
to $L$ as we like: however small a positive number we name, from some point
on every running total is within that distance of $L$. For
$\tfrac12+\tfrac14+\tfrac18+\dots$ the running total after $N$ numbers is
$1-\tfrac1{2^N}$ (Section~\ref{sec:intro}), and $\tfrac1{2^N}$ gets below any
positive number we name, since $2^N$ doubles at each step and so passes every
number. The sum is \emph{finite} when it adds up to some number, and
\emph{infinite} when its running totals pass every number.

\paragraph{Halving lists.} Putting $1$ in front of
$\tfrac12+\tfrac14+\dots=1$ gives $1+\tfrac12+\tfrac14+\dots=2$, and
multiplying every number by the same amount multiplies every running total,
and so the sum, by that amount. So a list in which each number is half the one
before adds up to twice its first number: for example
$\tfrac14+\tfrac18+\tfrac1{16}+\dots=\tfrac12$.

\paragraph{Running totals that only go up and stay below a fixed number add up
to a number.} Suppose the running totals only go up and never pass some fixed
number. Then they add up to
a number, and the decimals show why. Here are the running totals for every
exponent:
\begin{center}
\small
\begin{tabular}{c|c||c|c}
\papertablehead
fractions added & running total & fractions added & running total\\
$1$ & $1.0000\dots$ & $6$ & $1.5909\dots$\\
$2$ & $1.3333\dots$ & $7$ & $1.5988\dots$\\
$3$ & $1.4761\dots$ & $8$ & $1.6027\dots$\\
$4$ & $1.5428\dots$ & $9$ & $1.6047\dots$\\
$5$ & $1.5751\dots$ & $10$ & $1.6057\dots$
\end{tabular}
\end{center}
The whole-number part can only go up, and it stays below the fixed number, so
it changes only finitely often and then stays fixed; here it is $1$ from the
start. After that, the first digit after the point can only go up, and it
cannot pass $9$, so it also changes at most nine times and then stays fixed;
here it reads $0,3,4,5,5,5,5,6,6,6$. The same holds for the second digit, the
third, and every later one. The digits they settle on make one number, here
$1.6066951\dots$; the running totals never pass it, and once the first $d$
digits have settled they are within $\tfrac1{10^d}$ of it. So they add up to
that number.

\paragraph{Comparison.} If each number in one list is at most the matching
number in a second list, then each running total of the first is at most the
matching running total of the second. So if the second sum is finite, the
first is finite and no bigger; if moreover one number of the first list is
strictly smaller than its partner, the first sum is strictly smaller, by at
least the difference of those two numbers. Leaving numbers out of a list of
positive numbers only lowers its running totals.

\paragraph{The part after a cutoff.} For a finite sum, the numbers after a
cutoff add up to the whole sum minus the running total at the cutoff: for
$\tfrac12+\tfrac14+\dots=1$ cut after $\tfrac1{16}$, the rest is
$\tfrac1{32}+\tfrac1{64}+\dots=\tfrac1{16}=1-\tfrac{15}{16}$. Since the running
totals get as close to the sum as we like, the part after the cutoff gets as
small as we like. For a set of exponents $A$ and a cutoff $t$, the \emph{part
of $\sum_{n\in A}1/n$ beyond $t$} means $\sum_{n\in A,\,n>t}1/n$, the sum of
$1/n$ over the chosen exponents bigger than $t$; for the doubling exponents
and $t=4$ it is $\tfrac18+\tfrac1{16}+\dots=\tfrac14$.

\paragraph{Every choice of exponents gives a definite number.} Compare our
fractions with the halving list $1,\tfrac12,\tfrac14,\tfrac18,\dots$: exponent
by exponent, $1\le1$, $\tfrac13\le\tfrac12$, $\tfrac17\le\tfrac14$,
$\tfrac1{15}\le\tfrac18$, and in general $\tfrac1{2^n-1}\le\tfrac1{2^{n-1}}$,
because $2^n-1=2^{n-1}+(2^{n-1}-1)$ is at least $2^{n-1}$. The halving list adds
up to $2$, so for every set of exponents $A$ the running totals of
$\sum_{n\in A}1/(2^n-1)$ only go up and never pass $2$. So this sum is a
definite number $S_A$, at most $2$.

\subsection{The fraction test}\label{app:test}

\begin{lemma}[the fraction test]\label{lem:grid}
Let $x=p/q$, where $p$ and $q$ are whole numbers and $q\ge1$. If $m$ and $j$
are whole numbers and $mx-j$ is positive, then $mx-j\ge\tfrac1q$.
\end{lemma}

\begin{proof}
$mx-j=\tfrac{mp-jq}q$. The top is a whole number minus a whole number, and it
is positive because $mx-j$ is. So it is a positive whole number, at least $1$,
and $mx-j\ge\tfrac1q$.
\end{proof}

So to prove that a number $x$ is irrational it is enough to find, for every
positive integer $q$, whole numbers $m$ and $j$ (which may depend on $q$) with
$0<mx-j<\tfrac1q$: if $x$ were $p/q$, the lemma would give $mx-j\ge\tfrac1q$
for that $q$. The value $0$ must be excluded: $5\times\tfrac25-2=0$ is below
every $\tfrac1q$.

\subsection{Whole parts, fractional parts and \texorpdfstring{$R(k)$}{R(k)}}\label{app:split}

Dividing a whole number $a$ by a positive integer $b$ gives $a=sb+r$ with
whole numbers $s$ and $r$ and $0\le r<b$, so $\tfrac ab=s+\tfrac rb$. The
\emph{whole part} of $\tfrac ab$ is $s$ and its \emph{fractional part} is
$\tfrac rb$, which is at least $0$ and less than $1$. Only one such pair
exists: if $sb+r=s'b+r'$ with $r$ and $r'$ both between $0$ and $b-1$, then
$r-r'=(s'-s)b$ is a multiple of $b$ strictly between $-b$ and $b$, so it is
$0$. We write $k\rmod n$ for the remainder when $k$ is divided by $n$, so
$8\rmod3=2$. Adding a whole number $w$ to $\tfrac ab$ adds $w$ to the whole part
and leaves the fractional part unchanged.

Fix a set of exponents $A$ with sum $S=S_A$. For each chosen exponent $n$ and
each positive integer $k$, write $\lft_n(k)$ for the fractional part of
$\tfrac{2^k-1}{2^n-1}$, and define
\[
 R(k)=\sum_{n\in A}\lft_n(k).
\]
It adds the separate fractional parts, so it can be bigger than $1$. For the three
exponents $2,3,6$ and $k=5$, the multiplier is $31$, and
$\tfrac{31}3=10+\tfrac13$, $\tfrac{31}7=4+\tfrac37$,
$\tfrac{31}{63}=0+\tfrac{31}{63}$, so the fractional parts add up to
\[
 \frac13+\frac37+\frac{31}{63}=\frac{21+27+31}{63}=\frac{79}{63}=1+\frac{16}{63},
\]
while the whole product $31\times\bigl(\tfrac13+\tfrac17+\tfrac1{63}\bigr)=31\times\tfrac{31}{63}=\tfrac{961}{63}=15+\tfrac{16}{63}$
has fractional part $\tfrac{16}{63}$. The proof always uses the sum of the separate
fractional parts.

\begin{lemma}\label{lem:split}
For every positive integer $k$, the number $(2^k-1)S-R(k)$ is a whole number
$J(k)$. If $A$ is infinite, then $R(k)>0$.
\end{lemma}

\begin{proof}
An exponent $n>k$ has $2^n-1>2^k-1$, so its multiplied fraction is less than
$1$ and has whole part $0$. So only the finitely many chosen exponents up to
$k$ can have whole parts other than $0$, and the whole parts add up to a whole
number $J(k)$. Each fractional part
is at least $0$ and at most the multiplied fraction it comes from, and the
multiplied fractions add up to $(2^k-1)S$, so by comparison $R(k)$ is a definite
number. For the chosen exponents up to some $L\ge k$, splitting each of these
finitely many multiplied fractions gives
\[
 (2^k-1)\sum_{\substack{n\in A\\n\le L}}\frac1{2^n-1}
 =J(k)+\sum_{\substack{n\in A\\n\le L}}\lft_n(k).
\]
As $L$ grows, $J(k)$ does not change, and the two sums get as close as we like
to $S$ and to $R(k)$. Hence $(2^k-1)S=J(k)+R(k)$. If $A$ is infinite, it has a
chosen exponent $n>k$, because only $k$ positive integers are at most $k$; that
fraction's fractional part is all of $\tfrac{2^k-1}{2^n-1}$, which is positive, and
no fractional part is negative. So $R(k)>0$.
\end{proof}

\begin{corollary}\label{cor:atleast}
Let $A$ be infinite. If $S=p/q$ with whole numbers $p$ and $q\ge1$, then
$R(k)\ge\tfrac1q$ for every positive integer $k$.
\end{corollary}

\begin{proof}
$R(k)=(2^k-1)S-J(k)$ is positive, so the fraction test with $m=2^k-1$ and
$j=J(k)$ gives $R(k)\ge\tfrac1q$.
\end{proof}

Being infinite gives $R(k)>0$ for each $k$ separately, through a chosen
exponent beyond $k$; it gives no single positive number that all the $R(k)$
stay above (for the doubling exponents they get as small as we like). The
bound $\tfrac1q$ comes only from supposing $S=p/q$.

\begin{lemma}\label{lem:fracpart}
Let $k$ and $n$ be positive integers and $r=k\rmod n$. The fractional part of
$\tfrac{2^k-1}{2^n-1}$ is
\[
 \lft_n(k)=\frac{2^r-1}{2^n-1}.
\]
It is $0$ exactly when $n$ divides $k$, and it repeats every $n$ steps of $k$.
\end{lemma}

\begin{proof}
For every whole number $j$, $2^{j+n}-1=2^j(2^n-1)+(2^j-1)$, so
$\tfrac{2^{j+n}-1}{2^n-1}=2^j+\tfrac{2^j-1}{2^n-1}$: adding $n$ to the exponent
adds the whole number $2^j$ and leaves the fractional part unchanged. Write $k=sn+r$.
Applying this $s$ times, the fractional part at $k$ equals the fractional part at $r$. Since
$0\le r<n$, the fraction $\tfrac{2^r-1}{2^n-1}$ is at least $0$ and less than
$1$, so it is its own fractional part. It is $0$ exactly when $2^r=1$, that is, when
$r=0$, which says that $n$ divides $k$.
\end{proof}

\subsection{The doubling exponents}\label{app:doubling}

\begin{theorem}\label{thm:doubling}
The sum $\tfrac13+\tfrac1{15}+\tfrac1{255}+\tfrac1{65535}+\dots$ of the
fractions in exponents $2,4,8,16,\dots$ is irrational.
\end{theorem}

\begin{proof}
Let $k$ be one of the exponents $2,4,8,\dots$. Every chosen exponent up to $k$
divides $k$, so its fractional part is $0$ (Lemma~\ref{lem:fracpart}). Every later
chosen exponent $n$ is bigger than $k$, so its fractional part is the whole
multiplied fraction, and since $2^n-1\ge2^{n-1}$,
\[
 0<\frac{2^k-1}{2^n-1}<\frac{2^k}{2^{n-1}}=\frac2{2^{n-k}} .
\]
The gaps $n-k$ are different whole numbers, the smallest being $2k-k=k$.
Allowing every gap $e\ge k$ only adds positive numbers, so by comparison
\[
 0<R(k)<\sum_{e=k}^{\infty}\frac2{2^e}=\frac4{2^k}.
\]
If the sum were $p/q$, Corollary~\ref{cor:atleast} would give
$R(k)\ge\tfrac1q$ for every $k$. But the chosen exponents $k$ grow without
end, so some $k$ has $2^k>4q$, and then $R(k)<\tfrac1q$. This is a
contradiction.
\end{proof}

For example $R(8)=\tfrac1{257}+\tfrac1{16843009}+\dots\approx0.0039$, below the
bound $\tfrac1{64}$, and $R(16)\approx0.000015$, below $\tfrac1{16384}$.

\subsection{Averages over a full cycle}\label{app:cycle}

\begin{lemma}[average over a cycle]\label{lem:cycle}
Let $n$ and $Q$ be positive integers, let $g$ be their highest common factor
(the biggest whole number dividing both), and let $c=n/g$. As $m$ runs through
$1,2,\dots,c$, the remainders $mQ\rmod n$ are the numbers $0,g,2g,\dots,n-g$,
each exactly once, in some order, and after that they repeat with period $c$.
So for every positive multiple $M$ of $c$,
\[
 \frac1M\sum_{m=1}^{M}\lft_n(mQ)=\frac{g}{n(2^g-1)}-\frac1{2^n-1}<\frac1n .
\]
\end{lemma}

\begin{proof}
\emph{Period $c$.} Two values $m_1<m_2$ give the same remainder exactly when
$n$ divides $(m_2-m_1)Q$. Since $cQ=n\times\tfrac Qg$ and $\tfrac Qg$ is a
whole number, $n$ divides $cQ$, so moving $c$ steps on never changes the
remainder.

\emph{No shorter period.} Let $d$ be the smallest positive integer for which
$n$ divides $dQ$; so $d\le c$. Divide $c$ by $d$: $c=sd+r$ with $0\le r<d$.
Both $cQ$ and $sdQ$ are multiples of $n$, so $rQ$ is too, and since $r<d$, the
choice of $d$ gives $r=0$. So $d$ divides $c$, and hence $d$ divides $n=cg$.
Now $dQ=nf$ for a whole number $f$, so $Q=\tfrac nd\,f$: the whole number
$\tfrac nd$ divides both $n$ and $Q$, and therefore $\tfrac nd\le g$, that is,
$d\ge c$. So $d=c$.

\emph{Which remainders occur.} If two of $m=1,\dots,c$ gave the same
remainder, $n$ would divide $(m_2-m_1)Q$ with $0<m_2-m_1<c=d$, which is
impossible. So the first $c$ values give $c$ different remainders, each a
multiple of $g$ (because $g$ divides $mQ$ and $n$). There are exactly $c$
multiples of $g$ from $0$ to $n-g$, so each occurs once.

\emph{The average.} By Lemma~\ref{lem:fracpart} the fractional parts over one cycle
are $\tfrac{2^{ig}-1}{2^n-1}$ for $i=0,1,\dots,c-1$. Add $\tfrac1{2^n-1}$ to
each. With $G=1+2^g+2^{2g}+\dots+2^{(c-1)g}$, we have
$2^gG-G=2^{cg}-1=2^n-1$, so these add up to
\[
 \frac G{2^n-1}=\frac1{2^g-1}.
\]
Divide by the $c=\tfrac ng$ values and take away the amount added to each: the
average is $\tfrac g{n(2^g-1)}-\tfrac1{2^n-1}$. A whole number of cycles has
the same average. Finally $2^g-1=1+2+\dots+2^{g-1}\ge g$, so
$\tfrac g{n(2^g-1)}\le\tfrac1n$, and taking away the positive amount
$\tfrac1{2^n-1}$ leaves less than $\tfrac1n$.
\end{proof}

For $n=6$ and $Q=4$: $g=2$, $c=3$, and the average is
$\tfrac2{6\times3}-\tfrac1{63}=\tfrac19-\tfrac1{63}=\tfrac2{21}$, as in
Section~\ref{sec:general}. For $n=3$ and $Q=1$ it is
$\tfrac13-\tfrac17=\tfrac4{21}$.

\subsection{The bound for unfinished cycles}\label{app:crude}

\begin{lemma}\label{lem:crude}
For all positive integers $n$, $Q$ and $M$,
\[
 \frac1M\sum_{m=1}^{M}\lft_n(mQ)\le\frac{Q+2}{n} .
\]
\end{lemma}

\begin{proof}
\emph{One fractional part.} Fix $N\ge1$ and let $r=N\rmod n$. The distances $s\ge1$
from $N$ up to the later multiples of $n$ are $n-r,\,2n-r,\,3n-r,\dots$. Call
$X=\tfrac1{2^{n-r}}+\tfrac1{2^{2n-r}}+\tfrac1{2^{3n-r}}+\dots$, which is finite
by comparison with a halving list. Multiplying $X$ by $2^n$ gives $2^r+X$, so
$(2^n-1)X=2^r$ and $X=\tfrac{2^r}{2^n-1}$. Hence
\[
 \lft_n(N)=\frac{2^r-1}{2^n-1}<\frac{2^r}{2^n-1}
 =\frac1{2^{n-r}}+\frac1{2^{2n-r}}+\frac1{2^{3n-r}}+\dots
 =\sum_{\substack{s\ge1\\n\text{ divides }N+s}}\frac1{2^s} .
\]
For example $\lft_6(4)=\tfrac{15}{63}<\tfrac{16}{63}=\tfrac1{2^2}+\tfrac1{2^8}+\tfrac1{2^{14}}+\dots$,
the distances from $4$ to $6,12,18,\dots$ being $2,8,14,\dots$.

\emph{Cut off at $H$.} Fix a positive integer $H$. The terms with $s>H$ add up
to at most $\sum_{s>H}2^{-s}=2^{-H}$, so for each $m=1,\dots,M$,
\[
 \lft_n(mQ)<\sum_{\substack{1\le s\le H\\n\text{ divides }mQ+s}}\frac1{2^s}+\frac1{2^H}.
\]

\emph{Count by distance.} Add these $M$ inequalities. The right side now has
finitely many terms, so we may collect them by distance. For each $s$, let
$C_s$ be the number of $m$ for which $n$ divides $mQ+s$. The numbers
$Q+s,2Q+s,\dots,MQ+s$ are different and none is bigger than $MQ+s$, and there
are at most $\tfrac{MQ+s}n$ multiples of $n$ from $1$ to $MQ+s$; so
$C_s\le\tfrac{MQ+s}n$. Hence
\[
 \sum_{m=1}^{M}\lft_n(mQ)<\frac{MQ}n\sum_{s=1}^{H}\frac1{2^s}+\frac1n\sum_{s=1}^{H}\frac s{2^s}+\frac M{2^H}.
\]
The first sum is $1-\tfrac1{2^H}<1$. For the second, $U_H=\sum_{s=1}^{H}\tfrac s{2^s}$,
take away half of it from itself: the terms pair off as
$\tfrac s{2^s}-\tfrac{s-1}{2^s}=\tfrac1{2^s}$, so
$U_H-\tfrac{U_H}2=1-\tfrac1{2^H}-\tfrac H{2^{H+1}}$ and
$U_H=2-\tfrac{H+2}{2^H}<2$. So
$\sum_{m=1}^{M}\lft_n(mQ)<\tfrac{MQ+2}n+\tfrac M{2^H}$.

\emph{Let $H$ grow.} The left side does not depend on $H$, and
$\tfrac M{2^H}$ gets below any positive number we name, so the left side is at
most $\tfrac{MQ+2}n$. Dividing by $M\ge1$ gives at most
$\tfrac{Q+2/M}n\le\tfrac{Q+2}n$.
\end{proof}

For $n=6$, $Q=4$ and $M=2$, the actual average is
$\tfrac{15+3}{2\times63}=\tfrac17$ and the bound is $\tfrac{4+2}6=1$: weak, but
once $Q$ is fixed it is the same multiple of $\tfrac1n$ for every $M$.

\subsection{The proof of Theorem~\ref{thm:main}}\label{app:proof}

Let $A$ be infinite with $\sum_{n\in A}1/n$ finite, and suppose $S=p/q$ with
whole numbers $p$ and $q\ge1$. By Corollary~\ref{cor:atleast},
\begin{equation}\label{eq:atleast}
 R(k)\ge\frac1q\qquad\text{for every positive integer }k .
\end{equation}

\emph{The four choices.} By the part after a cutoff (Appendix~\ref{app:sums}),
choose $t$ with $\sum_{n\in A,\,n>t}\tfrac1n<\tfrac1{2q}$. Let $Q$ be the
product of the chosen exponents up to $t$, or $1$ if there are none; each of
them divides every $mQ$, so leaves $0$ there (Lemma~\ref{lem:fracpart}). With
$Q$ fixed, choose $T\ge t$ with $(Q+2)\sum_{n\in A,\,n>T}\tfrac1n<\tfrac1{2q}$.
Let $M$ be the product of the cycle lengths $\tfrac n{g}$ (with $g$ the highest
common factor of $n$ and $Q$) of the finitely many chosen exponents with
$t<n\le T$, or $1$ if there are none.

\emph{Adding by rows.} We need
\[
 \frac1M\sum_{m=1}^{M}R(mQ)=\sum_{n\in A}\frac1M\sum_{m=1}^{M}\lft_n(mQ).
\]
Keep only the rows $n\le H$; then both sides add the same finitely many
entries and agree. Every entry in a row $n$ is at most
$\tfrac{2^{MQ}-1}{2^n-1}\le(2^{MQ}-1)\tfrac2{2^n}$, so the rows $n>H$ add at
most $E_H=(2^{MQ}-1)\tfrac2{2^H}$ to each side. With $Q$ and $M$ fixed, $E_H$
gets below any positive number we name as $H$ grows, so the two sides are
equal.

\emph{The three groups.} Rows $n\le t$ contribute $0$. Each middle row
$t<n\le T$ finishes whole cycles, so by Lemma~\ref{lem:cycle} its average is
less than $\tfrac1n$, and together they contribute at most
$\sum_{n\in A,\,t<n\le T}\tfrac1n\le\sum_{n\in A,\,n>t}\tfrac1n<\tfrac1{2q}$
(``at most'' because the middle group may be empty). By
Lemma~\ref{lem:crude}, the far rows contribute at most
$(Q+2)\sum_{n\in A,\,n>T}\tfrac1n<\tfrac1{2q}$. So the average of
$R(Q),\dots,R(MQ)$ is less than $\tfrac1q$, and at least one of them is less
than $\tfrac1q$, contradicting~\eqref{eq:atleast}. Since $q$ was arbitrary,
$S$ is irrational.\qed

\subsection{The four choices with actual numbers}\label{app:numbers}

Take the doubling exponents and $q=5$, so the middle and far groups each get
$\tfrac1{2q}=\tfrac1{10}$. Beyond a doubling exponent $2^h$, the part of
$\sum_{n\in A}1/n$ is $\tfrac1{2^{h+1}}+\tfrac1{2^{h+2}}+\dots=\tfrac1{2^h}$.
\begin{itemize}
\item $t=16$: the part beyond it is $\tfrac1{16}<\tfrac1{10}$ ($t=8$ fails,
 since $\tfrac18>\tfrac1{10}$).
\item $Q=2\times4\times8\times16=1024$.
\item $T=16384$: $\tfrac{Q+2}{16384}=\tfrac{1026}{16384}<\tfrac1{10}$, since
 $10260<16384$ ($T=8192$ fails, since $10260>8192$).
\item The middle exponents are $32,64,\dots,16384$, with cycle lengths
 $1,1,1,1,1,1,2,4,8,16$ (each is $n$ divided by the smaller of $n$ and
 $1024$), so $M=2\times4\times8\times16=1024$.
\end{itemize}
The middle rows contribute less than
$\tfrac1{32}+\tfrac1{64}+\dots+\tfrac1{16384}=\tfrac{1023}{16384}$ and the far
rows at most $\tfrac{1026}{16384}$, so the average of $R(1024m)$ over
$m=1,\dots,1024$ is less than $\tfrac{2049}{16384}$, which is less than
$\tfrac15$ because $5\times2049=10245<16384$. So some $R(k)$ is below
$\tfrac15$, which rules out $S=p/5$ (in fact every bottom up to $7$, since
$7\times2049=14343<16384$). The largest $k$ used is $1024\times1024=1048576$;
the multiplier $2^{1048576}-1$ is never written out, because every step uses
only the bounds.

\subsection{Other bases}\label{app:bases}

Replace $2$ by a whole number $b\ge2$. For $b=3$ the list is
$\tfrac12,\tfrac18,\tfrac1{26},\tfrac1{80},\dots$, and exponents $2,4,8,\dots$
give $\tfrac18+\tfrac1{80}+\tfrac1{6560}+\dots$. Each step above carries over:
$b^n-1\ge2^{n-1}$, so the comparisons hold; the same calculation as in
Lemma~\ref{lem:fracpart} gives the fractional part $\tfrac{b^r-1}{b^n-1}$ with
$r=k\rmod n$; Lemma~\ref{lem:cycle} gives the average
$\tfrac g{n(b^g-1)}-\tfrac1{b^n-1}<\tfrac1n$, since $b^g-1\ge g$; and in
Lemma~\ref{lem:crude}, with $\tfrac1b$ in place of $\tfrac12$, the first sum is
at most $\tfrac1{b-1}\le1$ and the second at most $\tfrac b{(b-1)^2}=\tfrac1{b-1}+\tfrac1{(b-1)^2}\le2$,
so the bound $\tfrac{Q+2}n$ still holds. So for every $b\ge2$,
$\sum_{n\in A}1/(b^n-1)$ is irrational whenever $A$ is infinite and
$\sum_{n\in A}1/n$ is finite \cite[Theorem~A.1]{shortpaper}.

\subsection{The Lean statement}\label{app:lean}

\begin{center}
\begin{minipage}{0.9\linewidth}
\small\ttfamily
theorem irrational\_erdosSupportSeries\_of\_summable\_reciprocal\\
\hspace*{2em}(b : $\mathbb{N}$) (A : Set $\mathbb{N}$) (hb : 2 $\le$ b) (hA : A.Infinite)\\
\hspace*{2em}(hsum : Summable (reciprocalSupportTerm A)) :\\
\hspace*{2em}Irrational (erdosSupportSeries b A)
\end{minipage}
\end{center}
Line by line: $b$ is a whole number and \texttt{hb} assumes $b\ge2$; $A$ is a
set of whole numbers and \texttt{hA} assumes it is infinite;
\texttt{reciprocalSupportTerm A} is the list whose entry at $n$ is $1/n$ when
$n$ is in $A$ and $0$ otherwise, and \texttt{hsum} assumes its sum is finite;
the conclusion says that \texttt{erdosSupportSeries b A}, the sum of
$1/(b^n-1)$ over $A$, is irrational. Lean's whole numbers start at $0$, and in
Lean dividing by $0$ gives $0$, so the exponent $0$ counts as $0$ in both sums;
removing it changes neither sum and leaves an infinite set infinite. So for
$b=2$ the statement says the same as Theorem~\ref{thm:main}. It is the formal
version of \cite[Theorem~A.1]{shortpaper}.

The Lean kernel, a small program, has checked every step of the proof.
Comparator checked that the proof proves exactly a separately written formal
statement, one of eight theorems restated in the challenge
\texttt{E257\_47}, using only the axioms \texttt{propext},
\texttt{Quot.sound} and \texttt{Classical.choice}. The replay run 35935225572
passed, and its
\href{https://github.com/wcook04/plectis-erdos/blob/0064fc9b2a2693de97aebc8abec26090c3cd1095/evidence/comparator/replay-35935225572/receipt-E257_47.json}{Comparator receipt}
records acceptance by Lean's kernel and by a second, independent kernel
checker. At base $2$ the Lean proof follows \cite[Appendix~A]{shortpaper},
which takes limits of averages; it reaches the other bases by comparing each
base-$b$ fractional part with twice the base-$2$ one, a different route from
Appendix~\ref{app:bases}. The finite average of Appendix~\ref{app:proof} is
this paper's rewriting, checked by hand.

\section{Questions and exercises}\label{sec:exercises}

\begin{readercheck}\label{chk:hypotheses}
Which of $\{2,4\}$, $\{2,4,8,16,\dots\}$ and $\{2,4,6,8,\dots\}$ does
Theorem~\ref{thm:main} apply to?
\end{readercheck}

\begin{readercheck}\label{chk:positive}
If $A$ is infinite and $k$ cancels every chosen exponent up to $k$, why is
$R(k)$ still not $0$?
\end{readercheck}

\begin{readercheck}\label{chk:partial}
For the exponent $6$ and $Q=4$ the full cycle averages $\tfrac2{21}<\tfrac16$. Is
the first value alone also below $\tfrac16$?
\end{readercheck}

\begin{readercheck}\label{chk:order}
In Appendix~\ref{app:numbers}, replace $Q=1024$ by $Q'=4096$ and keep
$T=16384$. Does the far bound still give less than $\tfrac1{10}$? What must
be chosen again?
\end{readercheck}

\begin{exercise}\label{ex:1}
Write $\tfrac13+\tfrac17+\tfrac1{63}$ as a single fraction in lowest terms.
\end{exercise}

\begin{exercise}\label{ex:2}
List the fractional parts of $\tfrac{2^k-1}{15}$ for $k=1,\dots,8$, and check that
their average over one cycle is $\tfrac14-\tfrac1{15}$.
\end{exercise}

\begin{exercise}\label{ex:3}
Use Lemma~\ref{lem:fracpart} to show that $2^{12}-1$ is divisible by $7$ and
by $63$, and find its remainder on division by $2^5-1=31$.
\end{exercise}

\begin{exercise}\label{ex:4}
Adapt the proof of Theorem~\ref{thm:doubling} to the exponents
$3,9,27,81,\dots$, using $k=3^j$.
\end{exercise}

\begin{exercise}\label{ex:5}
Show that $5\sqrt2-7=(\sqrt2-1)^3$, and that multiplying a number
$a\sqrt2-c$ by $3-2\sqrt2$ gives another number of that form. Use the test of
Section~\ref{sec:test} to show that $\sqrt2$ is irrational.
\end{exercise}

\begin{exercise}\label{ex:6}
Let $a_1=2$ and let each later exponent be one more than the product of the
earlier ones, giving $2,3,7,43,1807,\dots$, and write $P_j=a_1a_2\cdots a_j$.
(a) Show that $\tfrac1{a_1}+\dots+\tfrac1{a_j}=1-\tfrac1{P_j}$. (b) Show that
with $k=P_j$ the first $j$ exponents leave $0$ but the exponent $P_j+1$ leaves at
least $\tfrac37$, so $R(P_j)\ge\tfrac37$ for every $j$.
\end{exercise}

\section{Answers}\label{app:answers}
\begingroup
\setlength{\parskip}{0.7ex plus 0.2ex}
\setlength{\parindent}{0pt}

\emph{Question~\ref{chk:hypotheses}.} Only $\{2,4,8,16,\dots\}$: $\{2,4\}$ is
finite (its sum $\tfrac25$ is a fraction), and for $\{2,4,6,\dots\}$ the sum
$\tfrac12+\tfrac14+\tfrac16+\dots$ is half of $1+\tfrac12+\tfrac13+\dots$,
which is infinite.

\emph{Question~\ref{chk:positive}.} $A$ is infinite, so some chosen exponent $n$
is bigger than $k$; that fraction times $2^k-1$ is positive and less than $1$,
so all of it is fractional part.

\emph{Question~\ref{chk:partial}.} No: the first value leaves
$\tfrac{15}{63}=\tfrac5{21}=\tfrac{10}{42}$, more than $\tfrac16=\tfrac7{42}$.
The full-cycle bound says nothing about part of a cycle, which is why
Lemma~\ref{lem:crude} is needed.

\emph{Question~\ref{chk:order}.} No:
$\tfrac{4098}{16384}=\tfrac{2049}{8192}$ is more than $\tfrac1{10}$, since
$20490>8192$. $T$ must be chosen again, for instance $T'=65536$, where
$\tfrac{4098}{65536}=\tfrac{2049}{32768}<\tfrac1{10}$, and then $M$ again. This
is why $Q$ is fixed before $T$.

\emph{Exercise~\ref{ex:1}.} $\tfrac{21}{63}+\tfrac9{63}+\tfrac1{63}=\tfrac{31}{63}$,
already in lowest terms, since $31$ is prime and does not divide $63$.

\emph{Exercise~\ref{ex:2}.} The remainders of $k=1,\dots,8$ divided by $4$ are
$1,2,3,0,1,2,3,0$, so the fractional parts are
$\tfrac1{15},\tfrac3{15},\tfrac7{15},0,\tfrac1{15},\tfrac3{15},\tfrac7{15},0$.
One cycle adds up to $\tfrac{11}{15}$, so the average is
$\tfrac{11}{60}=\tfrac14-\tfrac1{15}$.

\emph{Exercise~\ref{ex:3}.} $3$ and $6$ divide $12$, so $7$ and $63$ divide
$2^{12}-1=4095$. The remainder of $12$ divided by $5$ is $2$, so the remainder
on division by $31$ is $2^2-1=3$: indeed $4095=31\times132+3$.

\emph{Exercise~\ref{ex:4}.} For $k=3^j$, exponents $3^i$ with $i\le j$ divide
$k$ and leave nothing. Later exponents $n=3^i$ leave less than
$\tfrac2{2^{n-k}}$, and the gaps $3^i-3^j$ are different and at least
$2\times3^j$, so $R(3^j)<\sum_{e\ge2\times3^j}\tfrac2{2^e}=\tfrac4{2^{2\times3^j}}$,
which is below $\tfrac1q$ for large $j$.

\emph{Exercise~\ref{ex:5}.} $(\sqrt2-1)^2=3-2\sqrt2$ and
$(\sqrt2-1)(3-2\sqrt2)=5\sqrt2-7$. Also
$(a\sqrt2-c)(3-2\sqrt2)=(3a+2c)\sqrt2-(4a+3c)$. So every odd power
$(\sqrt2-1)^{2i+1}$ is $m\sqrt2-j$ for whole numbers $m$ and $j$, and it is
positive. Since $2.8^2=7.84<8$, we have $3-2\sqrt2<\tfrac15$, and
$\sqrt2-1<\tfrac12$, so $(\sqrt2-1)^{2i+1}<\tfrac12\times\tfrac1{5^i}$, which is
below $\tfrac1q$ once $5^i>q$. By the test, $\sqrt2$ is irrational.

\emph{Exercise~\ref{ex:6}.} (a) With $P_0=1$, $a_j=P_{j-1}+1$ and
$P_j=P_{j-1}a_j$, so
$\tfrac1{P_{j-1}}-\tfrac1{P_j}=\tfrac{a_j-1}{P_j}=\tfrac1{a_j}$, and adding
these for $1,\dots,j$ leaves $1-\tfrac1{P_j}$. (b) The first $j$ exponents
divide $P_j$. The exponent $P_j+1$ is bigger than $k=P_j$, so it leaves
$\tfrac{2^{P_j}-1}{2^{P_j+1}-1}$, which is at least $\tfrac37$ because
$7(2^{P_j}-1)-3(2^{P_j+1}-1)=2^{P_j}-4\ge0$. So cancelling longer and longer
early batches still leaves $R(P_j)\ge\tfrac37$, which is why the proof
averages.

\endgroup

\begin{thebibliography}{9}
\bibitem{erdos1948} P.~Erd\H{o}s, \emph{On arithmetical properties of Lambert
 series}, J.~Indian Math.\ Soc.\ 12 (1948), 63--66.
\bibitem{erdos1968} P.~Erd\H{o}s,
 \href{https://users.renyi.hu/~p_erdos/1969-09.pdf}{\emph{On the irrationality of
 certain series}}, Math.\ Student 36 (1968), 222--226 (issued 1969).
\bibitem{erdosproblems257} T.~F.~Bloom, \emph{Erd\H{o}s Problem \#257},
 \url{https://www.erdosproblems.com/257}; the statement is quoted in the Formal
 Conjectures file \texttt{FormalConjectures/ErdosProblems/257.lean}.
\bibitem{taoteravainen2025} T.~Tao and J.~Ter\"av\"ainen,
 \href{https://arxiv.org/abs/2512.01739v2}{\emph{Quantitative correlations and
 some problems on prime factors of consecutive integers}}, arXiv:2512.01739v2
 (submitted December 2025, revised April 2026).
\bibitem{shortpaper} W.~Cook,
 \href{https://github.com/wcook04/plectis-erdos/blob/0064fc9b2a2693de97aebc8abec26090c3cd1095/paper/257/erdos-257-mersenne-support-subseries.pdf}{\emph{Irrationality
 criteria for Lambert subseries}}, 30 September 2026.
\bibitem{longrecord} W.~Cook,
 \href{https://github.com/wcook04/plectis-erdos/blob/0064fc9b2a2693de97aebc8abec26090c3cd1095/paper/257/erdos257-mersenne-reasoning-surface.pdf}{\emph{Reciprocal
 Mersenne subseries: support criteria, finite arithmetic and conditional orbit
 problems}}, 30 September 2026.
\end{thebibliography}

\end{document}
